\documentclass[10pt,a4paper]{amsart}
\usepackage[margin=19mm]{geometry}
\usepackage{amsmath,amssymb,mathtools}
\usepackage[scr=rsfs,cal=euler]{mathalfa}
\usepackage{xfrac}
\usepackage[unicode,hidelinks,bookmarks=false]{hyperref}
\newcommand{\ie}{i.e.\ }
\newcommand{\cf}{cf.\ }
\newcommand{\resp}{respectively\ }
\newcommand{\Subsection}{\subsection}
\providecommand{\nobreakdash}{\nobreak\hbox{-}}
\setlength{\emergencystretch}{2em}
\multlinegap0pt
\usepackage{amssymb}
\usepackage{breqn}
\usepackage{enumerate}
\usepackage{bm}
\usepackage{bbm}
\usepackage{mathtools}
\usepackage{xcolor}
\usepackage{dsfont}
\newtheorem{theorem}{Theorem}
\newtheorem{claim}{Claim}
\newcommand\id{\mathbbm{1}}	
\newcommand{\on}[1]{
	\operatorname{#1}
}
\title{M\"obius function is strongly orthogonal to polynomial phases over $\mathbb{F}_p[t]$}
\author{Luka Mili\'cevi\'c, \v{Z}arko Ran{\dj}elovi\'c}
\date{Source: arXiv:2512.17633v1 (CC BY 4.0)}
\begin{document}
\maketitle

\noindent\emph{Source and licence.} M\"obius function is strongly orthogonal to polynomial phases over $\mathbb{F}_p[t]$. Version arXiv:2512.17633v1 (CC BY 4.0), distributed by arXiv under the Creative Commons Attribution 4.0 International licence, http://creativecommons.org/licenses/by/4.0/, as recorded in the arXiv licence metadata for this exact version and shown on the abstract page https://arxiv.org/abs/2512.17633v1. Full licence notice: \texttt{licenses/arxiv-2512.17633-CC-BY-4.0.txt}.\par\smallskip

\noindent\textit{Editorial scope.} Counted unit: the article's theorem biased\_ml\_form\_approx (source0.tex lines 677--681). The article's complete original proof of it is delivered with it (source0.tex lines 683--743, one \texttt{proof} environment of 61 lines). No mathematics is written, repaired or re-derived: the statement and the proof are the article's own lines, in the article's own notation, and no retained line is substituted or re-flowed.  No further source-local context is needed: every cross-reference of every retained line resolves inside the two delivered spans.  Nothing was omitted from any retained span, so the package carries no omission record.  No retained line contains a citation, so the wrapper carries no bibliography and no bibliography entry.  No retained line contains a figure environment, an image-inclusion command or drawing code, so no image is delivered and the package declares no asset dependency.  Editorial formatting only, in addition: an AMS \texttt{amsart} preamble that restates the article's own declarations and macros used by the retained lines, namely the environments \texttt{theorem}, \texttt{claim} on separate counters, together with the standard packages those lines need.  The retained environments print in this document as Theorem 1, Claim 1 in order of appearance, whereas the article prints the same items with the numbering of its own full document.  The complete native source of the article as distributed in the arXiv e-print for this exact version is bundled unchanged as \texttt{sources/191331/source0.tex}.
\begin{theorem}\label{biased_ml_form_approx}
    Let $\alpha : U_{[k]} \to \mathbb{F}_p$ be a multilinear form with $\on{bias}\alpha \geq c$. Let $\varepsilon > 0$. Then there exist
    $m \leq 2^{25}\varepsilon^{-14}c^{-14}$, coefficients $c_1, \dots, c_m \in \mathbb{C}$ such that $\sum_{i \in [m]} |c_i| \leq c^{-1}$, and multiaffine forms $\lambda_1, \dots, \lambda_m : U_{[k]} \to \mathbb{F}_p$ with vanishing multilinear part such that 
    \[\Big\| \chi \circ \alpha - \sum_{i \in [m]} c_i \chi \circ \lambda_i\Big\|_{L^2} \leq \varepsilon.\]
\end{theorem}

\begin{proof}
    Let $A: U_{[k-1]} \to U_k$ be the multilinear map such that $A(x_{[k-1]}) \cdot x_k = \alpha(x_{[k]})$. Write $Z = \{x_{[k-1]} \in U_{[k-1]} : A(x_{[k-1]}) = 0\}$, which has density at least $c$ by assumptions.\\

    Let $t_{[k-1]} \in U_{[k-1]}$ be arbitrary. Define the multiaffine map $L_{t_{[k-1]}} : U_{[k-1]} \to U_k$ by
    \[L_{t_{[k-1]}}(x_{[k-1]}) = \sum_{I \subsetneq [k-1]} (-1)^{k - |I|} A(x_I, t_{[k-1] \setminus I}).\]
    We remark that the multilinear part of $L_{t_{[k-1]}}$ vanishes since $I = [k-1]$ is not included in the sum.

    \begin{claim}\label{relationLandA}
        Whenever $x_{[k-1]} \in t_{[k-1]} + Z = \{(t+z)_{[k-1]} : z_{[k-1]} \in Z\}$, we have
        \[A(x_{[k-1]}) = L_{t_{[k-1]}}(x_{[k-1]}).\]
    \end{claim}

    \noindent\textbf{Remark.} The notation $(t+z)_{[k-1]}$ is a shorthand for the sequence $t_1 + z_1, \dots, t_{k-1} + z_{k-1}$ and $t_{[k-1]} + Z$ is simply a translate of $Z$ by $t_{[k-1]}$.\\

    \begin{proof}
        Take $z_{[k-1]} \in Z$ such that $x_i  = z_i + t_i$. Then
        \begin{align*}A(x_{[k-1]}) - L_{t_{[k-1]}}(x_{[k-1]}) = &A(x_{[k-1]})- \sum_{I \subsetneq [k-1]} (-1)^{k - |I|} A(x_I, t_{[k-1] \setminus I}) \\
        = &A(x_{[k-1]}) + \sum_{I \subsetneq [k-1]} (-1)^{k - 1 - |I|} A(x_I, t_{[k-1] \setminus I})\\
        = &\sum_{I \subseteq [k-1]} (-1)^{k - 1 - |I|} A(x_I, t_{[k-1] \setminus I})\\
        = &A(x_1 - t_1, \dots, x_{k-1} - t_{k-1}) = A(z_{[k-1]}) = 0.\qedhere\end{align*}
    \end{proof}

    \noindent\textbf{First approximation.} Let $m$ be a positive integer to be chosen later and let $t^{(1)}_{[k-1]}, \dots, t^{(m)}_{[k-1]} \in U_{[k-1]}$ be chosen uniformly and independently at random. By Hoeffding's inequality for $m$ independent random variables taking values in $[0,1]$, for each $x_{[k-1]} \in U_{[k-1]}$, the probability that 
    \begin{equation} \label{prbabilitybound} \Big||\{i \in [m]: x_{[k-1]} \in t^{(i)}_{[k-1]} + Z\}| - \frac{Z}{|U_{[k-1]}|}m\Big| \leq \frac{\varepsilon}{4} \frac{Z}{|U_{[k-1]}|}m \end{equation}
    is at least $1 - 2\exp(- 2^{-3}\varepsilon^2 c^2 m)$. Taking $m = 2^{7}\varepsilon^{-4} c^{-4}$ and using the elementary inequality $\exp(-M) = \frac{1}{\exp(M)} \leq \frac{1}{1 + M}\leq \frac{1}{M}$ for all $M > 0$, we get probability at least
    \[1 - 2\exp(- 2^{4}\varepsilon^{-2} c^{-2}) \geq 1 - \frac{c^2\varepsilon^2}{8}.\]

    \begin{claim}
        There exists a choice of $t^{(1)}_{[k-1]}, \dots, t^{(m)}_{[k-1]}$ such that the function 
        \[f(x_{[k]}) = \frac{|U_{[k-1]}|}{m |Z|} \sum_{i \in [m]} \id_{t^{(i)}_{[k-1]} + Z}(x_{[k-1]}) \chi\Big(L_{t^{(i)}_{[k-1]}}(x_{[k-1]}) \cdot x_k\Big)\]
        satisfies $\| \chi \circ \alpha - f\|_{L^2} \leq \varepsilon/2$.
    \end{claim}

    \begin{proof}
        Due to the probability bound above, we have a choice of $t^{(1)}_{[k-1]}, \dots, t^{(m)}_{[k-1]}$ such that~\eqref{prbabilitybound} holds on a set $X \subseteq U_{[k-1]}$ of size $|X| \geq (1 - \frac{c^2\varepsilon^2}{8})|U_{[k-1]}|$. Also, owing to Claim~\ref{relationLandA}, $\id_{t^{(i)}_{[k-1]} + Z}(x_{[k-1]}) \chi\Big(L_{t^{(i)}_{[k-1]}}(x_{[k-1]}) \cdot x_k\Big)$ is $\chi\Big(A(x_{[k-1]}) \cdot x_k\Big)$ when $x_{[k-1]} \in t^{(i)}_{[k-1]} + Z$, and 0 otherwise. Hence, we deduce that $|\chi \circ \alpha(x_{[k]}) - f(x_{[k]})| \leq \frac{|U_{[k-1]}|}{m |Z|} \frac{\varepsilon}{4} \frac{Z}{|U_{[k-1]}|}m = \frac{\varepsilon}{4}$ when $x_{[k-1]} \in X$, and $|\chi \circ \alpha(x_{[k]}) - f(x_{[k]})| \leq \frac{|U_{[k-1]}|}{|Z|} \leq c^{-1}$ for trivial reasons otherwise. Hence
        \begin{align*}\| \chi \circ \alpha - f\|_{L^2}^2 = &\frac{1}{|U_{[k-1]}|} \sum_{x_{[k-1]} \in U_{[k-1]}} |\chi \circ \alpha(x_{[k]}) - f(x_{[k]})|^2\\
        \leq &\frac{1}{|U_{[k-1]}|}\sum_{x_{[k-1]} \in X} \Big(\frac{\varepsilon}{4}\Big)^2 + \frac{1}{|U_{[k-1]}|}\sum_{x_{[k-1]} \in U_{[k-1]} \setminus X}c^{-2} \leq \varepsilon^2/4.\qedhere\end{align*}
    \end{proof}
    
    \vspace{\baselineskip}


    Let $t^{(1)}_{[k-1]}, \dots, t^{(m)}_{[k-1]}$ and $f$ be as in the claim.\\

    \vspace{\baselineskip}

    \noindent\textbf{Second approximation.} We may externally approximate $Z$ by a strictly-multilinear variety $V \subseteq U_{[k-1]}$ of codimension $s \leq 2\log_p (4\varepsilon^{-1} c^{-1} m)$ with $|V \setminus Z| \leq \Big(\frac{\varepsilon c}{4m}\Big)^2|U_{[k-1]}|$. Note that the $L^2$ norm of the function mapping $x_{[k]}$ to  
    \[\id_{t^{(i)}_{[k-1]} + V}(x_{[k-1]}) \chi\Big(L_{t^{(i)}_{[k-1]}}(x_{[k-1]}) \cdot x_k\Big) - \id_{t^{(i)}_{[k-1]} + Z}(x_{[k-1]}) \chi\Big(L_{t^{(i)}_{[k-1]}}(x_{[k-1]}) \cdot x_k\Big)\]
    is at most $\sqrt{|V \setminus Z|/|U_{[k-1]}|} \leq \frac{\varepsilon c}{4m}$.\\
    \indent Hence, defining 
    \[g(x_{[k]}) = \frac{|U_{[k-1]}|}{m |Z|} \sum_{i \in [m]} \id_{t^{(i)}_{[k-1]} + V}(x_{[k-1]}) \chi\Big(L_{t^{(i)}_{[k-1]}}(x_{[k-1]}) \cdot x_k\Big),\]
    by triangle inequality, we have $\| \chi \circ \alpha - g\|_{L^2} \leq \varepsilon$.\\

    \noindent\textbf{Final form of approximation.} Let $V = \beta^{(-1)}(0)$ for a multilinear map $\beta: U_{[k-1]} \to \mathbb{F}_p^s$. Then we may rewrite $g$ as 
    \begin{align*}g(x_{[k]}) = &\frac{|U_{[k-1]}|}{m |Z|} \sum_{i \in [m]} \id(\beta(x_1 - t^{(i)}_1,\dots,x_{k-1} - t^{(i)}_{k-1}) = 0) \chi\Big(L_{t^{(i)}_{[k-1]}}(x_{[k-1]}) \cdot x_k\Big)\\
    = & \frac{|U_{[k-1]}|}{m |Z|} \sum_{i \in [m]} p^{-s}\sum_{\tau \in \mathbb{F}_p^s} \chi\Big(\tau\cdot \beta(x_1 - t^{(i)}_1,\dots,x_{k-1} - t^{(i)}_{k-1})\Big) \chi\Big(L_{t^{(i)}_{[k-1]}}(x_{[k-1]}) \cdot x_k\Big)\\
    = &  \sum_{i \in [m], \tau \in \mathbb{F}_p^s} \frac{|U_{[k-1]}|}{p^s m |Z|} \chi\Big(\tau\cdot \beta(x_1 - t^{(i)}_1,\dots,x_{k-1} - t^{(i)}_{k-1}) + L_{t^{(i)}_{[k-1]}}(x_{[k-1]}) \cdot x_k\Big)
    \end{align*}
    which is the desired form of the approximation (we stress that we get phases of multiaffine forms with vanishing multilinear part). Note that the coefficients of the multiaffine phases are nonnegative reals, adding to $|U_{[k-1]}| / |Z| \leq c^{-1}$. The total number of terms in the approximation is 
    \[p^s m \leq (4\varepsilon^{-1} c^{-1} m)^2 m \leq 2^{25}\varepsilon^{-14}c^{-14}.\qedhere\]
\end{proof}

\end{document}
